\section*{Turing Machine Definitions}

A \kwub{turing machine (TM)} is specified by \iMbox{M=(Q, \Sigma, s, \delta)}
\begin{itemize}
    \vItem \kwb{machine states} \iMbox{Q} => finite set
    \vItem \kwb{tape symbols} \iMbox{\Sigma} => finite set,
    containing \stu{distinguished symbol} \iMbox{\left. \textvisiblespace \right.}, called \kwb{blank}
    \vItem \kwb{Initial state} \iMbox{s \in Q}
    \vItem \kwu{partial} \kwb{translation function}
    \iMbox{\delta \in(Q \times \Sigma) \rightharpoonup(Q \times \Sigma \times\{L, R\})}

\end{itemize}

\hSep

\kwub{TM configuration} => \iMbox{(q, w, u)}
\begin{enumerate}
    \vItem \kwb{current state} \iMbox{q \in Q}
    \vItem a \stu{finite (possibly-empty) string} \iMbox{w \in \Sigma^*} of \ul{tape symbols to the left} of tape head
    \vItem a \stu{finite (possibly-empty) string} \iMbox{u \in \Sigma^*} of \ul{tape symbols to the right} of tape head
    \vItem \textbf{NOTE:} \iMbox{\epsilon} is the \kwb{empty string}
    \vItem \kwb{initial configuration} \iMbox{(s, \epsilon, u)} => \stu{initial state} \iMbox{s} \& string of tape symbols \iMbox{u}
\end{enumerate}

\hSep

\iMbox{\mathrm{first}, \mathrm{last}: \Sigma^{*} \rightarrow \Sigma \times \Sigma^{*}} functions:
\begin{enumerate}
    \vItem \iMbox{\mathrm{first}(w) = \begin{cases}(a, v)                         & w=a v      \\
             ( \textvisiblespace, \epsilon) & w=\epsilon\end{cases}}
    \ \& \ \iMbox{\mathrm{last}(w) = \begin{cases}(a, v)                         & w=v a      \\
             ( \textvisiblespace, \epsilon) & w=\epsilon\end{cases}}
    \vItem These functions \ul{split off} the \stu{first and last symbols of a string}, splitting off \iMbox{\left. \textvisiblespace \right.}
    if the \stu{string is empty}
\end{enumerate}

\hSep

\mkImg[95pt]{image14}

\hSep

\kwub{TM normal form} => when \stu{configuration} \iMbox{M=(Q, \Sigma, s, \delta)}
has no computation step. \\
\textbf{NOTE:} this is \stu{exactly the case} when \iMbox{\delta(q,a) \uparrow} for \iMbox{\operatorname{first}(u) = (a, u^{\prime})}

\hSep

\kwub{TM computation} => finite \textit{(or infinite)} \stu{sequence} of \stu{configurations} \iMbox{c_{0}, c_{1}, c_{2}, \dots}
\begin{enumerate}
    \vItem \iMbox{c_{0}=\left( s, \epsilon, u \right)} is \stu{initial configuration}
    \vItem \iMbox{c_{i} \rightarrow_{M} c_{i+1}} holds for each \iMbox{i=0,1, \dots}
    \vItem \kwb{Non-halting Computation} => \stu{infinite sequence} \iMbox{c_{0}, c_{1}, \dots}
    \vItem \kwb{halting computation} => \stu{finite sequence} \iMbox{c_{0}, c_{1} \dots, c_{m}}
    where \iMbox{c_{m} = (q_m, w_m, u_m)} is a \stu{normal form}
\end{enumerate}

\stub{EXAMPLE: TM Computation}
\mkImg[90pt]{image18}

\hSep

\stub{EXAMPLE: TM Specification}
\mkImg{image15}

\hSep

\kwub{TM graphical representation}
\mkImg[90pt]{image16}

\stub{EXAMPLE: TM Diagram}
\mkImg{image17}

\hSep

\iMbox{f: \mathbb{N}^n \rightharpoonup \mathbb{N}} is \kwub{(Turing) computable function}
\stub[orange]{iff} there is TM \iMbox{M} such that:
\begin{enumerate}
    \vItem Starting \iMbox{M} from its \stu{initial state} with \stu{tape head} on
    \ul{leftmost \iMbox{0}} of \stu{tape coding} \iMbox{[x_1, \dots, x_n]}
    \vItem \iMbox{M} halts \stub[orange]{iff} \iMbox{f(x_1, \dots, x_n)\downarrow}, and \ul{in that case}:
    \begin{enumerate}
        \vItem \stu{Final tape} \ul{codes a list} whose \stu{first element} is \iMbox{y}
        \vItem where \iMbox{f(x_1, \dots, x_n) = y}
    \end{enumerate}
\end{enumerate}

\hSep

\subsection*{Turing Machine Tape Coding}

\stub{TM Tape Coding of Lists}:
\mkImg{image24}

\hSep

\stub[orange]{THEOREM:} The computation of a \stu{Turing machine} \iMbox{M} can be implemented by a \stu{register machine}

\mkImg[90pt]{image19}
\mkImg{image20}
\mkImg{image21}
\mkImg{image22}
\mkImg{image23}

\hSep

\stub[orange]{THEOREM:} A \stu{partial function} is \stu{Turing computable}
\stub[orange]{iff} it is \stu{RM computable}
\begin{enumerate}
    \vItem \stb{Proof (sketch).} => We already proved the forward direction. The backward direction
    would involve \stu{implementing an RM in terms of a TM}, which is doable in principle but loooooooong :)
\end{enumerate}

\kwub{Church-Turing Thesis} => Every \stu{algorithm [in intuitive sense]} can be
realized as a \stu{Turing machine}

\hSep
\vspace{2pt}

